\subsection{独立编码一致性}
问题一使用连续半开区间算法和离散时频单元倒排算法交叉检查；问题二对调整后的全部重复事件进行连续区间与离散资源格复核；问题三同时使用资源格集合装填和候选两两冲突两种编码；问题四对候选工作簿重新展开，并分别用连续区间、离散资源格和独立重建程序复核。四个问题的主要结果分别为 297 对、零冲突六撤销方案、141 台最大新增和三撤销可行见证，独立编码没有产生不一致。

\subsection{问题二的时间域敏感性}
\label{sec:q2-sensitivity}
统一时间域的右端点 643 来自完整数据展开。为检查该数据驱动边界是否决定撤销数量，分别考察 \(H=638,643,648\) 三个情景。每个情景都重新筛选越界候选，并对“总撤销数不超过 5”进行不可行检验，同时验证六撤销见证的边界和冲突状态。

\begin{table}[H]
\centering
\caption{不同时间域右端点下的最小撤销数}
\label{tab:q2-horizon}
\begin{tabular}{lrrl}
\toprule
时间域右端点 \(H\) & 候选动作数 & 最小撤销数 & 验证状态 \\
\midrule
638 & 4\,528 & 6 & 下界与见证均通过 \\
643 & 4\,582 & 6 & 基准方案通过 \\
648 & 4\,589 & 6 & 下界与宽域见证均通过 \\
\bottomrule
\end{tabular}
\end{table}

三个情景均得到最小撤销数 6，说明主目标对该范围内的时间边界扰动不敏感。但这项检验只支持“最小撤销数保持为 6”，不表示三个情景下的完整字典序末级平移量都已经证明最优。

表 \ref{tab:q2-horizon} 已经同时列出三个边界情景的候选动作数、最小撤销数和验证状态，足以支持该项敏感性结论，因此不再用图形重复表格。

\subsection{问题三的最大性复核}
问题三的集合装填模型和候选两两冲突模型分别得到 141 的上下界；增加 142 台阈值约束后返回不可行。该复核针对问题三的唯一主目标，且独立于问题二末级平移量的未闭合状态。

\subsection{问题四的候选覆盖与最小撤销层}
问题四先对有限动作全集做边界筛选，再检查支配替换的覆盖性。基准域内共有 6\,103 个合法候选动作，其中 5\,953 个为活动候选；支配消元保留 5\,533 个，删除的 420 个动作均能由同计划保留动作替换，候选覆盖核验通过。三撤销见证由连续区间、离散资源格和独立重建共同复核；对 $R=2$ 的六种类别分拆逐支建立阈值判定，六支均返回不可行。因此，$R_{Q4}^{\ast}=3$ 的最小撤销层在当前有限模型内形成上下界闭合。

\subsection{问题四条件候选与边界检查}
用于论文展示的条件候选含 147 个活动计划和 16\,356 个不同占用资源格；连续冲突、离散冲突和越界数均为 0，最大资源格占用数为 1，工作簿报告值与独立重算值一致。候选类别统计为 A 类 $0/20/0$、B 类 $0/39/1$、C 类 $5/83/2$（保留/调整/撤销），但 $R_B=1$、$M=142$ 和 $S_{10}^{(4)}=943$ 的相邻次级不可行证据尚未全部闭合，因此只作为条件候选呈现。进一步在 $H=638,643,648$ 下均找到通过验证的三撤销候选，支持可行上界对边界扰动稳定，不把两个敏感性域的最小值误写成已证明结论。

\subsection{跨问题接口检查}
问题三只读取问题二的冻结工作簿，不通过新增容量反向修改问题二；问题四则从问题一的原始 150 条计划独立起算，不读取 Q2/Q3 的批准或新增结果。合并验证同时检查资源边界、Q2--Q3 冲突、Q3 内部冲突、Q4 全量重复事件、重复资源格和输出工作簿行数，全部通过。因而四问结果形成“共同预处理—按问求解—按问复核”的单向证据链。
